


\chapter*{ABSTRACT}
\phantomsection
\addcontentsline{toc}{chapter}{ABSTRACT}


\bigskip

Large Language Models (LLMs) have exhibited remarkable capabilities in natural language understanding and code generation; however, their deployment in resource-constrained environments remains limited due to substantial computational and memory requirements. This study proposes an efficient methodology for Python code generation by fine-tuning the LLaMA-2 7B chat model utilizing the Quantized Low-Rank Adaptation (QLoRA) technique. A high-quality training dataset is curated and refined through a teacher-model-based filtering mechanism, while structured prompt design is employed to enhance logical consistency in generated outputs. Evaluation on the HumanEval benchmark demonstrates that the fine-tuned model improves Pass@1 from 0.1159 to 0.1280 and Pass@5 from 0.1890 to 0.2256 relative to the base model. Code quality assessment yields a BLEU score of 10.2703 and CodeBLEU of 0.2363, with quantization analysis indicating 16-bit BLEU/CodeBLEU of 5.8472/0.2003 versus 4-bit 5.7232/0.1945. Local deployment via Ollama enables efficient inference without reliance on cloud infrastructure, highlighting the effectiveness of parameter-efficient fine-tuning and optimized data curation in developing practical, high-performance code generation systems.
\vspace{1em}

\noindent\textit{\textbf{Keywords}}---LLMs, LLAMA2, Code Generation,PEFT, QLoRA, LoRA, Hugging Face, Ollama, Local Inference, Prompt Engineering.

% \chapter*{ABSTRACT}
% \onehalfspacing % Ensures consistent spacing as per campus standards
% Fine-tuning large language models requires intensive hardware resources, making full fine-tuning infeasible to train under low-resource constraints. Although some large language models perform better in Python code generation, they are not feasible for training on custom datasets. Addressing this gap requires efficient fine-tuning that reduces memory and computation overhead. This project aims to fine-tune Llama-2 7B for Python code generation by Parameter Efficient Fine-Tuning (PEFT) with Quantized Low Rank Adaptation (QLoRA), focusing on hardware-level constraints. The training dataset is constructed from multiple open-source repositories, including FlyTec, StaQC, Alpaca 18k and TACO. Preprocessing involves filtering non-Python samples, cleaning noisy code, and formatting the data into instruction–response pairs suitable for supervised fine-tuning. The fine-tuned model will be evaluated on unseen Python problems. Performance will be compared with the base model Llama-2 7B to assess the improvement achieved through fine-tuning. The fine-tuned model is expected to be capable of generating Python code solutions for common coding exercises, using a mixture of common programming problems and LeetCode problems as a benchmark on low-edge devices.

% \vspace{1cm}
% \noindent \textbf{Keywords:} PEFT, QLoRA, LoRA, OPT, Hugging Face.
